\section{Admissible figures and complete support maps}
\label{sec:figure-comparisons}

\subsection{Normality and its finite interval consequence}
We first prove Theorem~\ref{thm:flags}. Apply
Proposition~\ref{prop:figure} to the singleton blocks of a clique.
It says that admissibility of $F$ is equivalent to
\begin{equation}\label{eq:clique-test}
 z\in\cl F,\quad F\subseteq N(z)\cup N(w)
 \quad\Longrightarrow\quad z\perp w.
\end{equation}
A maximal clique has empty polar and hence closure $E$. Thus
admissibility of every clique gives the stated maximal-clique
condition. Conversely extend $F$ to a maximal clique $D$. Every
$d\in D\setminus F$ lies in $N(F)$ and is orthogonal to any
$z\in\cl F$. Therefore the premise in \eqref{eq:clique-test}
implies $D\subseteq N(z)\cup N(w)$. The case $z=w$ is impossible:
it would put $z$ in both $N(F)$ and $\cl F$. The maximal-clique
condition gives $z\perp w$, as required. This proves the equivalence,
including infinite cliques when maximal extensions are chosen.

For the finite interval assertion, give a point symbol equal left
and right endpoints and give $I_{ab}$ endpoints $a,b$. Define the
positional order by
\[
 Z<_{\mathrm p}W\quad\Longleftrightarrow\quad
 Z\ne W\ \hbox{ and the right endpoint of $Z$ is at most the
 left endpoint of $W$}.
\]
It is a strict partial order, and its comparable pairs are exactly
the orthogonal pairs. A finite orthogonal subset is therefore a
finite positional chain. Its endpoints form a chain
$t_0<\cdots<t_n$ in $L$. Insert every $P_{t_i}$ and each
$I_{t_it_{i+1}}$ to obtain an alternating chain $H$. It contains
the original subset: an original interval cannot contain another
endpoint strictly between its endpoints, since the corresponding
symbol would not be orthogonal to it.

For any symbol $Z$, let $T_Z$ consist of the members of $H$ not
orthogonal to $Z$. If nonempty, this is a contiguous block.
Indeed, a member between two incomparable members can be neither
below nor above $Z$, by transitivity. Every member before the block
is below $Z$, and every member after it is above $Z$.
Suppose two nonempty blocks $T_Z,T_W$ are disjoint, with $T_Z$
first. At a consecutive cut $h_i,h_{i+1}$ between them,
$Z<_{\mathrm p}h_{i+1}$ and $h_i<_{\mathrm p}W$.
The consecutive symbols are $P_t,I_{tu}$ or $I_{st},P_t$.
Thus the right endpoint of $Z$ is at most $t$, and the left
endpoint of $W$ is at least $t$. It follows that
$Z<_{\mathrm p}W$. The only possible failure of strictness would
be the same point symbol twice, which the disjoint nonempty
blocks exclude.

Now take $Z\in\cl H$ and $H\subseteq N(Z)\cup N(W)$.
The first trace is nonempty, and the two traces are disjoint.
If $T_W$ is empty, $W\in N(H)$ gives $Z\perp W$ directly.
Otherwise the preceding cut argument gives the same conclusion.
Equation~\eqref{eq:clique-test} proves admissibility of $H$.
For a subchain $F\subseteq H$ and $A\subseteq F$,
$\cl F\subseteq N(H\setminus F)$, whence
\[
 \cl F\cap N(A)
 \subseteq\cl H\cap N((H\setminus F)\cup A)
 =\cl(F\setminus A).
\]
The reverse inclusion is immediate. Empty subsets are included.
This proves the finite assertion in Theorem~\ref{thm:flags}.

\begin{counterexample}[Finite figure complements do not force orthomodularity]
\label{ex:eight-symbols}
Let $L=\{a,b,c,d\}$ with exactly
$a<c,a<d,b<c,b<d$. The eight-symbol orthoset $M(L)$ has twenty
supports and is normal, but its support lattice is not orthomodular.
\end{counterexample}
\begin{proof}
There are four point edges and, for each interval symbol, the two
edges joining it to its endpoints. The supports are the empty set,
the eight singletons, the six pairs of point symbols, the four
point neighborhoods, and the whole set.
Here $N(I_{xy})=\{P_x,P_y\}$, while intersecting the neighborhoods
of an upper and a lower endpoint gives their interval singleton.
Intersecting the two upper point neighborhoods gives
$\{P_a,P_b\}$, and dually for the two lower neighborhoods.
Further intersections give the point singletons. The displayed
family is closed under intersections and contains all point and
interval neighborhoods, so it is precisely the family of all
polars, establishing the count.

The closed sets $A=\{P_a\}\subset B=\{P_a,P_b\}$ satisfy
$B\cap N(A)=\varnothing$, contradicting the orthomodular law.
The maximal cliques are the four endpoint--interval triangles.
If such a triangle lies in $N(u)\cup N(v)$, its interval forces
one of $u,v$ to be an endpoint. Covering that endpoint forces the
other vertex to be its neighbor. Thus the maximal-clique condition
holds. The finite-flag result proves the same figure-complement
conclusion without the enumeration.
\end{proof}

\begin{counterexample}[An infinite flag across a cut]\label{ex:infinite-flag}
Let $L=\Q\amalg\{a,b\}$, with the usual rational order,
$a<q$ exactly for $q\ge0$, and $q<b$ exactly for $q<0$.
There are no further comparisons involving $a,b$.
The orthogonal family $F=\{P_q:q\in\Q\}$ is not admissible.
\end{counterexample}
\begin{proof}
Its polar is empty. A rational point fails at itself; $P_a$ misses
negative rationals and $P_b$ misses nonnegative rationals. An
interval with rational endpoints misses an interior rational.
An interval starting at $a$ misses negative rationals, and one
ending at $b$ misses nonnegative rationals. These exhaust the
possible symbols. Thus $\cl F=M(L)$.
Put $A=\{P_q:q\ge0\}$ and $B=F\setminus A$.
Then $P_a\in N(A)$ and $P_b\in N(B)$, but $P_a\not\perp P_b$.
Consequently $P_a\notin\cl B$ although
$P_a\in\cl F\cap N(A)$.
\end{proof}

\subsection{Which families detect the orthomodular law?}
\begin{proposition}\label{prop:orthogonal-families}
In a complete orthomodular lattice, every mutually orthogonal
family $(e_i)_{i\in I}$ of nonzero elements gives an injective
complete Boolean embedding into $[0,\bigvee_i e_i]$ by
$A\mapsto\bigvee_{i\in A}e_i$.
Conversely, the relative-complement identity for every two-element
orthogonal family of nonzero elements implies orthomodularity.
\end{proposition}
\begin{proof}
For orthogonal $u,v$, the dual orthomodular identity gives
$(u\vee v)\wedge u^\perp=v$. Apply this with
$u=\bigvee_{i\in A}e_i$ and $v=\bigvee_{i\notin A}e_i$.
These joins are orthogonal, even for an infinite family, by
completeness and reversal of order under complement.
The relative complement follows. Joins are preserved directly;
intersections follow by complementing unions inside the displayed
interval. An index in $A\setminus B$ contributes a nonzero element
below the first join and orthogonal to the second, so the map is
injective.

Conversely let $a\le b$. Use the two orthogonal elements
$u=b^\perp$ and $v=a$, omitting a zero term. The assumed identity
gives $(b^\perp\vee a)\wedge b=a$. Complementing gives the
orthomodular law for $b^\perp\le a^\perp$, hence for every
comparable pair. Counterexample~\ref{ex:eight-symbols} shows why
testing only the families supplied by figures is weaker.
\end{proof}

\begin{proposition}[Faithful refinement and regular opens]\label{prop:regular-open}
Let $\rho$ be a preorder on $E$ and define $x\perp y$ to mean that
there is no $z$ with $z\rho x$ and $z\rho y$. Then $\Sigma(E)$
is the complete Boolean algebra of regular opens of the lower
Alexandrov topology of $\rho$.
\end{proposition}
\begin{proof}
Write $O_A=\bigcup_{a\in A}\mathord\downarrow_\rho a$.
A point is orthogonal to $A$ exactly when its principal lower set
misses $O_A$. Therefore
\[
 N(A)=\operatorname{int}(E\setminus O_A),\qquad
 \cl A=\operatorname{int}\overline{O_A}.
\]
The latter is regular open, and every regular open $U$ is obtained
by taking $A=U$. Two disjoint regular opens are orthogonal because
a common lower refinement of points in them would lie in their
intersection. The Boolean argument in Section~2 applies.
Arbitrary intersections of regular opens are open in an Alexandrov
topology and remain regular by monotonicity of
$\operatorname{int}\overline{(-)}$. These are the complete meets;
joins are regular-open closures of unions and complement is
interior of the set-theoretic complement.
\end{proof}

\begin{proposition}[A sufficient graph class with a strict converse]
\label{prop:cographs}
Every finite cograph has an orthomodular support lattice. A
five-vertex comparability graph can have a Boolean support lattice
without being a cograph.
\end{proposition}
\begin{proof}
Define cographs recursively from a one-vertex graph by disjoint
union and graph join. For a graph join, taking polars acts
coordinatewise on the vertex parts, so its support lattice is the
product of the support lattices of those parts. For a disjoint
union, every proper nonzero closed set lies in one part and is
proper and closed there. Complements agree with the complement
in that part. Thus the support lattice is the horizontal sum,
identifying all zeros and all tops; a two-element factor adds no
proper element. Products and horizontal sums preserve
orthomodularity. For the latter, every comparable proper pair
lies in a single factor. Induction proves the first assertion.

For strictness use five points represented by
$\{1\},\{2\},\{3\},\{1,2\},\{2,3\}$, orthogonal when
disjoint. The three singleton points detect the union represented
by any subset of points. Its double polar consists exactly of
the five points whose representing subsets lie in that union.
All eight unions occur, so the support lattice is the Boolean
algebra on three atoms. Calling the points $a,b,c,x,y$ in the
displayed order, the graph has edges $ab,ac,bc,cx,ay$.
It is the comparability graph of $a<b<c$, $x<c$, $a<y$.
The path $x-c-a-y$ is induced, which a cograph cannot contain:
an induced four-vertex path is neither a nontrivial disjoint union nor a
nontrivial join, and the recursive construction is inherited by induced
subgraphs. This proves strictness.
\end{proof}

\subsection{Restriction to a set of points}
The following criterion identifies exactly when a smaller orthoset
determines every support. Let $L\subseteq M$ have the orthogonality
induced from an orthoset $M$. Suppose each $x\in M$ has a nonempty
set $U_x\subseteq L$, with $U_p=\{p\}$ for $p\in L$, such that
\begin{equation}\label{eq:lieu-detection}
 x\perp y\quad\Longleftrightarrow\quad U_x\perp U_y.
\end{equation}
Orthogonality on the right means every pair is orthogonal.

\begin{proposition}[Exact support restriction]\label{prop:lieu-supports}
Under these hypotheses, the maps
\[
 S\longmapsto S\cap L,\qquad
 T\longmapsto\widehat T=\{x\in M:U_x\subseteq T\}
\]
are inverse complete ortholattice isomorphisms
$\Sigma(M)\simeq\Sigma(L)$.
Consequently the orthomodular and Boolean criteria are exactly
the corresponding criteria on $L$. Every orthoset is a possible
point orthoset in this statement.
\end{proposition}
\begin{proof}
Put $U(A)=\bigcup_{x\in A}U_x$. Detection gives
\[
 N_M(A)=\widehat{N_L(U(A))},\qquad
 \widehat T\cap L=T.
\]
If $T$ is closed in $L$, then
$N_M(\widehat T)=\widehat{N_L(T)}$: one implication tests the
points $p\in T\subseteq\widehat T$, and the other uses detection
and $U_y\subseteq T$ for $y\in\widehat T$. Applying this twice
proves that $\widehat T$ is closed in $M$. Conversely a closed
support $S=N_M(A)$ is the hat of the closed polar $N_L(U(A))$.
The trace formula proves that these maps are inverse; the polar
formula proves that they preserve orthocomplements. An order
isomorphism preserves all joins and meets. Taking $M=L$ and
$U_x=\{x\}$ realizes any given orthoset, including the empty one.
\end{proof}

In a refinement presentation the points called \emph{lieux} are
the minimal refinement objects, and $U_x$ is the set of lieux
internally refining $x$. Minimality gives $U_p=\{p\}$.
Thus nonempty open lieu sets and the explicit detection
condition~\eqref{eq:lieu-detection} are the precise inputs of this
support comparison. No additional restriction on the point graph
follows from the general formalism: incidence and refinement can
both be identity, with the given orthogonality, so every object
is a lieu and each principal deployment is a singleton.

\paragraph{Faithfulness and lieu atoms.}
For a strong system with nonempty open-lieu sets and local detection,
the following conditions are equivalent: faithfulness; pairwise
disjunction of distinct lieux; and Booleanity of its support lattice
with distinct lieux generating distinct atoms.
Indeed, distinct lieux have no common internal refinement, so
faithfulness forces their disjunction. Conversely, if the lieu
graph is complete, local detection says that two strata are
nondisjoint exactly when their open-lieu sets intersect. A common
lieu is a common internal refinement; any common internal
refinement contains a lieu by nonemptiness and internal
transitivity. This is faithfulness.
The induced support lattice on a complete lieu graph is
$\mathcal P(L)$, and Proposition~\ref{prop:lieu-supports} gives
the corresponding statement for all strata. Conversely distinct
atoms of a Boolean ortholattice are orthogonal, so distinct
lieux must be disjoint.

If the lieu graph additionally has singleton closed supports,
Booleanity alone forces distinct lieux to be disjoint: two
different singleton atoms in a Boolean algebra are orthogonal.
Without that separation condition, two nondisjoint lieux may
generate the same Boolean atom, as the two-point example in
Section~2 shows. For a connected comparability poset $L$, the
nonisolated-core criterion in Theorem~\ref{thm:interval} reduces
to totality of $L$, with the singleton case included.

The separation issue also occurs with genuine interval strata.
Take $L=\mathbb Q\amalg\{c\}$, with $c$ incomparable to every
rational. Its nonisolated core is the chain $\mathbb Q$, so its
support lattice is Boolean. Every rational interval has a nonempty
open-lieu set by density, and all point open-lieu sets are
singletons. Local detection holds: two rational interval interiors
overlap exactly when they have a common rational point, which
witnesses failure of disjunction; the point--interval case is the
same membership test. Distinct rational points are disjoint.
The point $c$ is disjoint from no rational point, and every open
lieu set being tested is nonempty, so detection also holds for
pairs involving $c$. Nevertheless $P_c$ is nondisjoint from every
other point and has full support because its polar is empty.
Thus nonempty open-lieu sets and Booleanity do not replace the
stated separation condition.

\subsection{Objects and correspondences}
We prove Theorem~\ref{thm:correspondence} in full. On
$E_K=K\setminus\{0\}$ the relation $x\le Cy$ is symmetric and
irreflexive. For $A\subseteq E_K$ its polar is the nonzero
principal ideal below $C(\bigvee A)$, and its double polar is the
nonzero principal ideal below $\bigvee A$. Every principal ideal
occurs, and inclusion and complement give the asserted complete
ortholattice isomorphism. The empty orthoset represents the
one-element lattice.

The zero convention has a precise variant. Suppose a symmetric
relation permits self-orthogonal elements only when they are
orthogonal to every element, and let $R$ be their set. Every polar
contains $R$. Removing $R$ gives an ordinary orthoset, and
$S\mapsto S\setminus R$, with inverse adjoining $R$, identifies
the two complete support lattices. Thus retaining zero does not
silently change the irreflexive convention.

If $f$ preserves all joins, then
$\cl A=\bigvee_{a\in A}\cl\{a\}$ proves the correspondence
condition for $\alpha(x)=f(\cl\{x\})$. Conversely that condition
gives
\[
 \bigvee_{x\in\cl A}\alpha(x)=\bigvee_{a\in A}\alpha(a).
\]
It follows immediately that $f(S)=\bigvee_{x\in S}\alpha(x)$
preserves all joins. Taking $A=\{x\}$ and using
$x\in\cl\{x\}$ shows that these assignments are inverse.
Substitution gives the identity and composition formulas.
The object representation just proved makes the resulting
full and faithful functor essentially surjective.

Every complete join map has right adjoint
$g(T)=\bigvee\{S:f(S)\le T\}$. Hence
\[
 f(S)\le C_FT\quad\Longleftrightarrow\quad
 S\le g(C_FT)\quad\Longleftrightarrow\quad
 S\perp C_Eg(C_FT).
\]
The conjugate $C_EgC_F$ preserves joins since $g$ preserves meets.
The displayed characterization proves uniqueness, reverses
composition, and gives $(f^\dagger)^\dagger=f$.
Equivalently the correspondence condition says that each column
$\{x:\alpha(x)\perp y\}$ is closed: apply the condition to
that column for one direction, and test all elements of the
polar of $\bigvee_{a\in A}\alpha(a)$ for the other.

Finally, for a classical orthoset, where distinct points are
orthogonal, supports are all subsets. Complete join maps between
such powersets are arbitrary relations. A relation sending one
point to two is not a function, and a function identifying two
points need not preserve orthogonality. Both give the specified
correspondences. This proves the stated distinction of categories.
